\section{From formal series to entire functions}
\subsection{Finite exponential sums are the analytic realization}
The entire transfer function is defined as a finite root-weighted exponential sum.
\excerpt{EntireBridge}{entire_transfer}
The primitive root, residue index, and spectral parameter appear explicitly. The finite sum is entire because each exponential summand is entire. The crucial next step is a \code{HasSum} theorem identifying its Taylor coefficients with the formal transfer section.
\excerpt{EntireBridge}{entire_sum}
This theorem has positive order, a residue index in range, and a primitive root as hypotheses. It provides a convergent series at every complex argument. It also permits independence-of-root proofs: two finite sums obtained from allowed root choices have the same summing series, hence the same value by uniqueness of a sum.

The signed analytic realization additionally chooses $\rho$ with $\rho^N=-1$. The factor $\rho^{-j}$ compensates for the residue label. It cannot be removed without changing the normalization of the basis. The independent-root statements compare the resulting functions under exactly the root assumptions they require.

\subsection{Initial data and Wronskians}
The formal cycle specifies derivatives at zero, while the analytic bridge transfers them to iterated complex derivatives. For the normalized signed basis, the matrix of initial data is the identity. Constancy of the Wronskian then identifies the Wronskian with one for the stated $N>1$ range.
\excerpt{Wronskians}{wronskian}
The source does not conclude linear independence from an informal appeal to distinct exponentials alone. It factors the exponential jet matrix, handles its determinant, proves the Wronskian law, and identifies the normalized initial matrix. The ODE modules then describe the corresponding realization of the solution space.

\subsection{Companion matrices and finite algebras}
The signed shift matrix has a negative carry across the cyclic boundary. Its powers encode paths through the residue labels. This provides a companion realization of the signed differential cycle, explicit fundamental-matrix entries, addition and inversion identities, and a spectrum controlled by the chosen roots.

There are two compatible algebra models. In the cyclic quotient algebra the generator satisfies $x^N=1$. In the signed quotient it satisfies $x^N=-1$. Earlier modules construct the quotient algebras and their power bases, then prove formal exponential identities. Row 151 equips the finite-dimensional quotients with coefficient norm topologies and proves convergence of the actual exponential series. Consequently, the analytic identity in a finite algebra is more than equality of formal coefficient expressions.

\subsection{Why the final pointwise addition theorem matters}
A bivariate addition law for formal series is an algebraic identity among coefficients. It does not on its own prove equality of evaluated functions at arbitrary complex arguments. The completed source includes both. Row 144 obtains the actual transfer and signed addition formulas, retaining the signed carry factor forced by $\rho^N=-1$.

The useful reading order is: formal product and addition identities, the entire summation bridge, then the pointwise closure. Keeping these stages visible prevents a notation change from hiding the convergence theorem. It also makes the formal statements reusable in contexts where no analytic topology has been chosen.

\subsection{Growth order and indicator geometry}
The entire functions are finite exponential sums. Their exponents form regular polygon configurations in the complex plane. In directions with a unique exposed vertex, one exponential dominates. On tie directions, cancellation can obstruct a naive pointwise lower bound, so the exact indicator proof uses aligned radii in addition to uniform sector estimates.

Rows 110--115 prove sharp order, exponential type, indicator geometry, uniform dominance in sectors, and the exact directional indicator. These results are proved in the source. Imported standard fact 11.1 remains recorded separately and is not required for the checked order proofs. The guide therefore does not count that imported fact among the 68 completed results.
